\documentclass{../notatki}

\title{Optyka z fotoniką}

\begin{document}

\section{Wektorowe operatory różniczkowe}

Czas na powtórkę i wstęp z okropnej notacji fizycznej.
$$
\vec{f}: \mathbb{R}^3 \to \mathbb{R}^3
$$
$$
\vec{f}(x, y, z) = f_x \hat{i} + f_y \hat{j} + f_z \hat{k}
$$

\subsection{Gradient}

$$
\nabla \vec{f}(x, y, z) =
\hat{i} \frac{\partial f}{\partial x}(x, y,z) +
\hat{j} \frac{\partial f}{\partial y}(x, y, z) +
\hat{k} \frac{\partial f}{\partial z}(x, y, z)
\in \mathbb{R}^3
$$

\subsection{Dywergencja}

$$
\nabla \cdot \vec{f} = \frac{\partial f_x}{\partial x} + \frac{\partial
f_y}{\partial y} + \frac{\partial f_z}{\partial z}
\in \mathbb{R}
$$

Mierzy jak bardzo pole wektorowe rozbiega się w danym punkcie.

\subsection{Rotacja}

$$
\nabla \times \vec{f}(x, y, z) =
\begin{vmatrix}
  \hat{i} & \hat{j} & \hat{k} \\
  \frac{\partial}{\partial x} & \frac{\partial}{\partial y} &
  \frac{\partial}{\partial z} \\
  f_x(x, y, z) & f_y(x, y, z) & f_z(x, y, z)
\end{vmatrix}
\in \mathbb{R}^3
$$

Mierzy jak bardzo pole wektorowe kręci się dookoła danego punktu.

\subsection{Zależności}

$$
\nabla \cdot (\nabla \times \vec{f}) = 0
$$
$$
\nabla \times (\nabla f) = 0
$$
$$
\nabla \times (\nabla \times \vec{f}) = \nabla(\nabla \cdot \vec{f})
- \Delta \vec{f}
$$

\subsection{Laplasjan}

$$
\nabla \cdot (\nabla f(x, y, z)) \equiv \nabla ^2 f(x, y, z) =
\Delta f(x, y, z) =
\frac{\partial^2 f}{\partial^2 x}(x, y, z) +
\frac{\partial^2 f}{\partial^2 y}(x, y, z) +
\frac{\partial^2 f}{\partial^2 z}(x, y, z)
\in \mathbb{R}^3
$$

\subsection{Iloczyn wektorowy}

$$
\vec{c} = \vec{a} \times \vec{b} =
\begin{vmatrix}
  \hat{i} & \hat{j} & \hat{k} \\
  a_x & a_y & a_z \\
  b_x & b_y & b_z
\end{vmatrix}
$$

\subsection{Twierdzenie Gaussa-Ostrogradskiego}

$$
\int_V \nabla \cdot \vec{f} dx dy dz = \oint_A \vec{f} \cdot \vec{dA}
$$
Całka po objętości $V$ z dywergencji, jest równa całce z powierzchni otaczającej
daną objętość.

\subsection{Twierdzenie Stokesa}

$$
\oint_C \vec{f} \cdot \vec{dl} = \int_A (\nabla \times \vec{f}) \cdot \vec{dA}
$$
gdzie $C$ to dowolna krzywa, a A to powierzchnia dowolnego otwartego obszaru
ograniczonego krzywą C.

\section{Równania Maxwella}

\subsection{Prawo Faradaya}

$$
\oint_L \vec{E} \vec{dl} = - \frac{\partial \phi_B}{\partial t} = -
\frac{d}{dt} \int_S \vec{B} \vec{ds}
$$
gdzie $\vec{E}$ to natężenie pola elektrycznego, $L$ to dowolna zamknięta
krzywa, $\phi_B$ to strumień indukcji przez powierzchnię $S$ rozpinającą $L$,
a $\vec{B}$ to indukcja pola magnetycznego.

$$
\nabla \times \vec{E} = - \frac{\partial \vec{B}}{\partial t}
$$

\subsection{Prawo Amper'a}

$$
\oint_L \vec{B} \vec{dl} = \mu I + \mu \epsilon \frac{\partial
\phi_E}{\partial t} = \mu I + \mu \epsilon \frac{d
}{dt} \int_S \vec{E} \vec{ds}
$$
gdzie $I$ to całkowity prąd elektryczny przepływający przez dowolną powierzchnię
$S$ rozpinającą $L$, $\mu$ i $\epsilon$ to stałe.

$$
\nabla \times \vec{B} = \mu \vec{g} + \mu \epsilon \frac{\partial
\vec{E}}{\partial t}
$$
gdzie $\vec{g}$ to gęstość prądu elektrycznego.

\subsection{Prawo Gaussa}

$$
\epsilon \phi_E = \epsilon \oint_S \vec{E} \vec{ds} = q
$$
gdzie $\phi_E$ to strumień pola elektrycznego przez dowolną zamkniętą
powierzchnię $S$, a $q$ to całkowity ładunek zawarty wewnątrz tej powierzchni.

$$
\epsilon \nabla \cdot \vec{E} = \rho
$$
gdzie $\rho$ to gęstość ładunku elektrycznego.

\subsection{Prawo Gaussa dla magnetyzmu}

$$
\phi_B = \oint_S \vec{B} \vec{ds}  = 0
$$

$$
\nabla \times \vec{B} = 0
$$

\section{Wektor Poyntinga}

$$
\vec{S} = \frac{1}{\mu_0}(\vec{E} \times \vec{B})
$$

\section{Siła Lorentza}

$$
\vec{F} = q \vec{E} + q(\vec{v} \times B)
$$

\section{Funkcje falowe}

Funkcja falowa, ma postać ogólną:
$$
f(x, t) = A \sin(kx + \omega t)
$$
lub
$$
f(x, t) = Ae^{i(kz - \omega t)}
$$

$$
\omega = 2 \pi f \quad\left[\frac{\textrm{rad}}{\textrm{s}}\right]
$$
$$
f = 1/T \quad\left[\textrm{Hz}\right] \text{lub} \left[\textrm{s}^{-1}\right]
$$
$$
k = \frac{2 \pi}{\lambda} \quad\left[\frac{\textrm{rad}}{\textrm{m}}\right]
$$
$$
\lambda = \frac{v}{f} \quad[\textrm{m}]
$$
$$
v = f\lambda = \frac{\omega}{k} \quad\left[\frac{\textrm{m}}{\textrm{s}}\right]
$$

\section{Wzór Eulera}

$$
e^{ix} = \cos(x) + i \sin(x)
$$

Tutaj przydają się powszechne uproszczenia funkcji trygonometrycznych:
\begin{table}[h!]
  \centering
  \renewcommand{\arraystretch}{1.3}
  \begin{tabular}{|c|c|c|c|c|}
    \hline
    $\alpha$ & $30^\circ\ (\frac{\pi}{6})$ &
    $45^\circ\ (\frac{\pi}{4})$ & $60^\circ\ (\frac{\pi}{3})$ &
    $90^\circ\ (\frac{\pi}{2})$ \\
    \hline
    $\sin \alpha$   & $\frac{1}{2}$ & $\frac{\sqrt{2}}{2}$ &
    $\frac{\sqrt{3}}{2}$ & $1$ \\
    \hline
    $\cos \alpha$   & $\frac{\sqrt{3}}{2}$ & $\frac{\sqrt{2}}{2}$ &
    $\frac{1}{2}$ & $0$ \\
    \hline
    $\tan \alpha$   & $\frac{1}{\sqrt{3}}$ lub  $\frac{\sqrt{3}}{3}$
    & $1$ & $\sqrt{3}$ & nie istnieje \\
    \hline
    $\cot \alpha$   & $\sqrt{3}$ & $1$ & $\frac{1}{\sqrt{3}}$ lub
    $\frac{\sqrt{3}}{3}$ & $0$ \\
    \hline
  \end{tabular}
  \caption{Wartości funkcji trygonometrycznych dla wybranych kątów}
\end{table}

\section{Załamanie Światła}

\begin{figure}[h!]
  \centering
  \begin{tikzpicture}
    \draw[dotted, ->] (0,-2) -- (0,2) node[above] {$y$};
    \draw[dotted, ->] (-2,0) -- (2,0) node[right] {$x$};

    \draw[->, thick] (-1.5,1.5) node[above left] {$\vec{k}_I$} -- (0,0);
    \draw[->, thick] (0,0) -- (1.2,1.5) node[above right] {$\vec{k}_R$};
    \draw[->, thick] (0,0) -- (0.8,-1.5) node[below right] {$\vec{k}_T$};
  \end{tikzpicture}
  \caption{Monochromatyczna fala elektromagnetyczna padająca z
    ośrodka 1 ($\vec{k}_I$). Część fali ulega odbiciu ($\vec{k}_R$),
    a inna część
  załamaniu ($\vec{k}_T$)}
\end{figure}

Wszystkie fale mają tą samą częstość, stąd:
$$
\omega = k_Iv_1 = k_Rv_1 = k_Tv_2
$$
Co pozwala nam określić, że:
$$
k_I = k_R \quad k_R = k_T \frac{v_2}{v_1} = k_T \frac{n_1}{n_2}
$$
gdzie $n_1$ to współczynnik załamania ośrodka, $v_1$ to prędkość
światła w ośrodku.

Kąt padania i odbicia są równe:
$$
\theta_I = \theta_R
$$

\subsection{Prawo Snella}

Przy światłu odbijającemu się od wody:
$$
n_1 \sin(\theta_1) = n_2 \sin(\theta_2)
$$
gdzie $n_1$ to współczynnik załamania ośrodka z którego światło pada, $\theta_1$
to kąt padania, a $\theta_2$ to kąt odbicia. Prawo Snella pozwla określić
kąty załamania i odbicia, oraz współczynniki załamania, lecz nie pozwala:
\begin{itemize}
  \item określić jaka część fali zostanie odbita
  \item jak się zmienia amplituda czy natężenie fali
  \item jak się zmienia polaryzacja
\end{itemize}

\begin{examplebox}
  Oblicz kąt załamania i kąt odbicia dla kąta padania 30 stopni i
  współczynników załamania $n_1=1$ i $n_2=3$
  $$
  \sin(\theta_1) = \frac{n_2}{n_1} \sin(\theta_2) \Rightarrow
  \sin(30^\circ) = \frac{1}{2} \sin(\theta_2)
  $$
  $$
  \sin(\theta_2) = 2\sin(30^\circ) \Rightarrow \sin(\theta_2) = 1
  \Rightarrow \theta_2 = 90^\circ
  $$
\end{examplebox}

\subsection{Wzory Fresnela}

Dla światła spolaryzowanego prostopadle $\perp$ do płaszczyzny padania:
$$
r_\perp = \frac{n_1 \cos(\theta_1) - n_2 \cos(\theta_2)}{n_1
\cos(\theta_1) + n_2 \cos(\theta_2)} \quad t_\perp = \frac{2 n_1
\cos(\theta_1)}{n_1
\cos(\theta_1) + n_2 \cos(\theta_2)}
$$
gdzie $t_\perp$ to współczynnik amplitudowy przepuszczania, a
$r_\perp$ to współczynnik amplitudowy odbicia.

Dla światłą spolaryzowanego równolegle do płaszczyzny padania:
$$
r_\parallel = \frac{n_2 \cos(\theta_1) - n_1 \cos(\theta_2)}{n_2
\cos(\theta_1) + n_1 \cos(\theta_2)} \quad t_\parallel = \frac{2 n_1
\cos(\theta_1)}{n_2 \cos(\theta_1) + n_1 \cos(\theta_2)}
$$

\begin{examplebox}
  Oblicz $t_\perp$ dla $\theta_1=60^\circ$ i $\theta_2=10^\circ$. Podaj
  przykłady $n_1$ i $n_2$ pasujące do tych kątów.
  $$
  t_\perp = \frac{2 n_1 \cos(60^\circ)}{n_1 \cos(60^\circ) + n_2
  \cos(10^\circ)} = \frac{n_1}{n_1 \frac{1}{2} + n_2 \cos(10^\circ)}
  $$
  $$
  n_1 \sin(60^\circ) = n_2 \sin(10^\circ) \Rightarrow n_2 =
  \frac{n_1 \sqrt{3}}{2 \sin(10^\circ)}
  $$
  $$
  t_\perp = \frac{n_1}{\frac{1}{2} \left(n_1 + \frac{n_1 \sqrt{3}
  \cos(10^\circ)}{\sin(10^\circ)} \right)} = \frac{2 n_1}{n_1 +
  \cot(10^\circ) n_1 \sqrt{3}} = \frac{2}{1 + \cot(10^\circ) \sqrt{3}}
  $$
\end{examplebox}

\subsection{Kąt Brewstera}

$$
\tan(\theta_P) = \frac{n_2}{n_1}
$$
gdzie $\theta_P$ to kąt Brewstera, lub inaczej kąt padania, pod którym światło
pada na granicę dwóch ośrodków i światło odbite jest całkowicie spolaryzowane.
Kąt istnieje dla $n_1 \neq n_2$, lecz musi występować polaryzacja równoległa.

W pewnym sensie, kąt Brewstera to $\theta_1$, jeśli spełnione są warunki.

\subsection{Kąt graniczny}

$$
\sin(\theta_C) = \frac{n_2}{n_1}
$$
gdzie $\theta_C$ to kąt graniczny, czyli kąt padania, pod jakim światło może
przejść z jednego ośrodka do drugiego. $n_1 > n_2$.

Kąt graniczny wychodzi z Prawa Snella, stąd niezależy od warunków polaryzacji.

Fala zanikająca, to taka fala, która pada na granice dwóch ośrodków
od strony większego współczynnika załamania i zanika,
ponieważ jej kąt padania jest większy od kąta granicznego.

\begin{examplebox}
  Na przykład, dla $n_1 = 1.5$ i $n_2 = 1.0$, mamy:
  $$
  \sin(\theta_C) = \frac{n_2}{n_1} = \frac{1.0}{1.5} = \frac{2}{3}
  \Rightarrow \theta_C \approx 42^\circ
  $$
  Zatem, dla $\theta_1 > \theta_C = 42^\circ$, światło zaniknie.
\end{examplebox}

\section{Model Drudego}

Ośrodki, w których funkcja dielektryczna zależy od częstotliwości nazywamy
materiałami dyspersyjnymi. Dla takich materiałów:
$$
\omega_p^2 = \frac{ne^2}{\varepsilon_0m} \quad \varepsilon(\omega) =
1 + i\frac{\omega_p^2}{\omega(\gamma - i\omega)}
$$
gdzie $\varepsilon(\omega)$ to całkowita przenikalność dielektryczna,
$\varepsilon_0$ to przenikalność dielektryczna w próżni,
$\omega_p$ to plazmonowa częstość własna zależna od materiału,
$\omega$ to częstość fali elektromagnetycznej. Jeśli zaniedbamy tłumienie:
$$
\varepsilon(\omega) = 1 - \frac{\omega_p^2}{\omega^2}
$$

Model Drudego, czyli wyrażenie na przewodność elektronów swobodnych to:
$$
\sigma = \frac{ne^2}{m(\gamma - i\omega)}
$$

Dla idealnego, niestraconego metalu (bez tłumienia, $\gamma=0$) i
dielektryka o przenikalności ($\varepsilon_d$):
$$
\omega = \frac{\omega_p}{\sqrt{1 + \varepsilon_d}}
$$

Model Drudego bez tłumienia nie przewiduje strat, więc aby obliczyć długość
propagacji fali elektromagnetycznej, trzeba uwzględnić straty.

\end{document}
